\section{Related work}
\label{sec:related}

\paragraph{Lie-group integrators}
Integrators that advance rotations in a local tangent chart and return to the group through the
exponential map go back to Crouch and Grossman~\cite{crouch1993numerical} and to the
Runge--Kutta--Munthe-Kaas construction~\cite{munthe1998lie,munthe1999high}; the surveys of
Iserles et al.~\cite{iserles2000lie} and of Celledoni, Marthinsen and Owren~\cite{celledoni2014lie}
and the monograph of Hairer, Lubich and Wanner~\cite{hairer2006geometric} give the general theory.
In multibody dynamics the same viewpoint underlies the Lie-group generalized-$\alpha$ schemes of
Br\"uls, Cardona and Arnold~\cite{bruls2010lie,bruls2012lie}, whose convergence for constrained
systems was analysed in~\cite{arnold2007convergence,arnold2015error}, BDF schemes on Lie
groups~\cite{wieloch2021bdf}, the geometrically exact beam and rigid-body formulations on
$SE(3)$~\cite{sonneville2014geometrically}, and singularity-free quaternion
integration~\cite{terze2016singularity,terze2020aircraft}. Bottasso and Borri~\cite{bottasso1998integrating}
discuss the interaction between the rotation parametrization and the integrator. These methods
establish how to evolve rotations geometrically; they are of order at most two, or, in the
Runge--Kutta--Munthe-Kaas case, are formulated for ordinary differential equations on the group
rather than for the constrained index-3 problem.

\paragraph{Constrained mechanical systems and index-3 DAEs}
Absolute-coordinate multibody systems are index-3 differential-algebraic equations
\cite{hairer1993solving,ascher1998computer,petzold1986numerical}. Their numerical order depends
not only on the integrator but on how the constraints and their differentiated forms are enforced:
index reduction with stabilization~\cite{gear1985automatic}, projection, and the review of
constraint-enforcement approaches by Bauchau and Laulusa~\cite{bauchau2008review} describe the
options and their failure modes. Runge--Kutta methods for index-2 and overdetermined index-3
formulations were analysed by Jay~\cite{jay1996symplectic,jay2006specialized,jay2007specialized};
the Lobatto IIIA--IIIB pairs of that work are the closest relatives of the present stage system in
that they enforce the position and velocity constraints at every stage. Lunk and
Simeon~\cite{lunk2006solving} and Arnold and Br\"uls~\cite{arnold2007convergence} treat the
Newmark and generalized-$\alpha$ families for constrained systems. The absolute-coordinate
formulations of Negrut and co-workers
\cite{taves2020exponential,kissel2021dwelling,kissel2022absolute,fang2023halfimplicit,kissel2024velocitypartitioning}
compare three rotation parametrizations and a half-implicit scheme on a public suite of four
mechanisms; that suite and its public code are used in Section~\ref{sec:numerical}. Reduced
(joint-coordinate) formulations, for which the constrained problem becomes an ordinary
differential equation, are classical~\cite{nikravesh1988computer,shabana2020dynamics,jain2011robot};
the present paper shows that a particular absolute-coordinate collocation step coincides with
Gauss collocation of such a reduced formulation without forming it.

\paragraph{Collocation and variational integrators}
Gauss collocation is the standard high-order implicit Runge--Kutta family: the $s$-stage method has
order $2s$, is symplectic and $B$-stable, and its stages are the values of a collocation polynomial
\cite{hairer1993nonstiff,butcher2016numerical,hairer2006geometric}. Variational and symplectic
integrators~\cite{marsden2001discrete,leyendecker2008variational,leok2012prolongation,bourabee2009hamilton}
derive the discrete equations from a discrete action and treat holonomic constraints at the
discrete level; Betsch and Steinmann~\cite{betsch2002conservation} obtain energy--momentum
conserving time finite elements for constrained systems. The present method is not variational: it
enforces the constraint hierarchy pointwise at the stages and leaves the conservation properties
to those of Gauss collocation on the reduced problem.

\paragraph{Friction and nonsmooth dynamics}
Regularized friction laws such as the Brown--McPhee model~\cite{browmcphee2016friction} replace the
set-valued Coulomb law by a smooth velocity-dependent force; Pennestr\`{\i} et
al.~\cite{pennestri2016review} review the alternatives. Nonsmooth time-stepping and complementarity
formulations~\cite{stewart1996implicit,anitescu1997formulating,tasora2011matrixfree} handle
unilateral contact and Coulomb friction without regularization but are of low order by design. The
problems in this paper are bilateral lower-pair mechanisms with regularized friction; the effect of
the regularization parameter on the observed order is quantified in Section~\ref{sec:numerical}.

\paragraph{Time finite elements}
Time finite elements for structural and rigid-body dynamics~\cite{borri1988rigid,hughes1988spacetime,hulbert1992time,kim2017higher,jog2016robust,bauchau2024time}
provide a different route to high order. The Lie-group time finite-element method of Chaturvedi,
Sandu and Sandu~\cite{chaturvedi2026tfe} extends this route to frictional index-3 multibody DAEs;
its Gauss--Lobatto family of degree $m$ has formal order $2m-1$, so that its $m=3$ member is a
fifth-order formula, although the source paper reports order reduction on the DAE problem. We do
not reimplement that method; the comparison with it in this paper is at the level of formal order.
